\section{Failure Analysis}
\label{sec:failures}

The failures are the most informative part of this study, and two of them
contradict what we expected before running it.

\subsection{Vector RAG lost the category it was expected to win}

T1 asks what one named company disclosed about one topic in its risk factors. The
answer is entirely inside a single document, retrieval is a ranking problem, and
similarity search should be at its best. It passed 3 of 15 attempts, all from one
question.

The mechanism is visible in the retrieved set. Asked what manufacturing and
supply-chain risks \emph{Moderna} disclosed, the top-8 chunks came from Entegris,
Hayward Holdings, Open Text, Arrowhead Pharmaceuticals and NVIDIA. Every one of
those chunks genuinely discusses supply-chain risk. The retriever did what it was
asked. The problem is that ``belongs to Moderna'' is not a property the embedding
encodes, so the company name in the question cannot outweigh topical similarity
across 508{,}714 chunks drawn from 5{,}210 filers.

The decisive detail is that \textbf{this question passed on a single month of
filings}. The failure appeared when the corpus grew. It is therefore a
scale-induced limitation of entity filtering by similarity, not a tuning error, and
it is the strongest argument in this paper for structural retrieval --- and one we
did not anticipate.

We are careful about the boundary of this claim. This experiment cannot fully
separate ``embeddings cannot filter by entity'' from ``this retriever
configuration cannot'', because the 1{,}400-character chunks exceed the encoder's
256-token window and $k$ was fixed at 8. A hybrid retriever with a metadata filter
on CIK would very likely fix T1 --- which is precisely the point, since a metadata
filter on a stable identifier \emph{is} structural retrieval, arriving by another
route.

\subsection{Capability and safety moved in opposite directions}

The benchmark's most uncomfortable result is not a score but a direction of
travel. While calibrating the harness we found that both stores normalise names
by stripping corporate suffixes and that neither schema description said so
(\Cref{sec:harness}). Documenting it identically for every arm --- one sentence
of schema prose, no change to retrieval code, questions or scoring --- moved
three of the four categories.

It also made the strongest baseline \emph{more dangerous}. With the schema fully
described, text-to-SQL stopped refusing multi-hop questions and started answering
them, and its confident falsehoods rose accordingly to the six reported here.
Better schema knowledge did not make it more careful. It converted honest
abstentions into assertive answers, and on this corpus an assertive multi-hop
answer keyed on a person's name is exactly how three different people called
Lozano become one.

This is worth separating from the accuracy comparison, because the two do not
move together. Giving a text-to-SQL system a better description of its own schema
is an unambiguous capability improvement and it made the system less safe on the
questions where identity is ambiguous. GraphRAG's falsehood count was zero
throughout, not because it is a better writer of queries --- it needed its repair
allowance far more often --- but because identity is resolved on an identifier
before any text is read, so the wrong people are excluded structurally rather
than by the model choosing correctly.

\subsection{GraphRAG lost aggregation, as predicted}

Text-to-SQL scored 15 of 15 on T2 and GraphRAG 6 of 15. Unlike the previous two,
this outcome was expected --- T2 exists in the benchmark precisely because a
traversal should add nothing to a \texttt{GROUP BY} --- but the manner of the loss
is instructive.

Asked for the most common jurisdiction of incorporation among all disclosed
subsidiaries, the generated Cypher grouped by company rather than by jurisdiction,
returned per-company counts, and the model reported Delaware with 2{,}415
subsidiaries. Delaware is the right answer; 2{,}415 is one company's count rather
than the 62{,}504 across the corpus. Right entity, wrong aggregate. The answer
also attached 190 accession numbers to a question that needs none.

For counting over typed columns the graph adds no information and the extra hop
adds a place to go wrong, at 24 times the token cost.

\subsection{Every falsehood was a surname collision}

Six of the 180 attempts stated something verifiably false. All six were
text-to-SQL, all six on the multi-hop category, and all six are the same error,
spread across four of its five questions (\texttt{T3-1}, \texttt{T3-2}, \texttt{T3-3}, \texttt{T3-4}). Asked about \emph{Monica Lozano},
the model reported companies including Mondelez and Ternium; those belong to
Mariano and Santiago Lozano, different people who share a surname. Asked about
\emph{William Giles}, it reported ten companies.

The SQL was not malformed and the model was fluent and confident. It matched on
name because name is what the question supplied and what the text contains. This
is the exact failure that resolving identity to an identifier before reading any
text prevents structurally rather than probabilistically: GraphRAG passed 14 of
15 attempts on this category and made no false statement anywhere in the
benchmark, because the other Lozanos never entered its result set.

One caution about single-attempt evaluation follows. Of the three attempts
text-to-SQL made on the featured question the outcomes were not identical, and an
evaluation that drew only one of them could easily have recorded an honest
refusal and missed the failure entirely. Four of the five multi-hop questions
produced at least one falsehood and at least one non-falsehood, so on this
category a single draw per question is close to a coin toss over whether the most
important failure mode in the paper is observed at all.

\subsection{The verdict is model-sensitive}

Because all arms route through one gateway (\Cref{sec:arch}), we probed the same
GraphRAG pipeline on other models. \Cref{tab:crossmodel} reports the featured
multi-hop question: the benchmarked model passes, Claude Haiku 4.5 fails because
its stage-1 Cypher collapsed to a single company and missed two others, and Amazon
Nova 2 Lite refuses.

This is a probe and not a measurement --- one question, one attempt each, no
re-scoring of the full set --- and we draw only the negative conclusion it
supports: a single-model benchmark bounds what can be claimed. The architecture
constrains what evidence reaches the model, but the model still has to write a
correct traversal, and not all of them do.
